\section*{第 \ref{ch:b3:behavioural-ecology} 章——行为生态学}

\begin{solution}{exo:b3:behavioural-ecology:1}
机制：舞蹈者将飞行中测得的食物与太阳之间的夹角，
换算为巢脾上相对重力的角度，并
由飞行途中的光流估算距离，并将其编码为摆尾段的
时长。发育：基本是先天的——从未见过舞蹈的蜜蜂，
仍能跳出正确的舞蹈，只有少量后天习得的调整。功能：
招募巢友前往有利可图的食物源，提高蜂群的摄食量。
历史：由离巢取食的仪式化意向运动演化而来；
侏儒蜜蜂仍在水平表面上跳舞，
直接指向食物源，这是祖先形式。
\end{solution}

\begin{solution}{exo:b3:behavioural-ecology:2}
\omterm{def:b3:behavioural-ecology:tinbergen}{固定动作模式}是由简单线索释放、一旦开始便会完整进行的刻板行为序列，
该线索即\omterm{def:b3:behavioural-ecology:tinbergen}{信号刺激}：鹅滚卵的行为，
由卵形物体释放；海鸥雏鸟啄食的行为，
由类似喙形物上的红色斑点释放。一根带三条红色条纹的木棍，
被啄的次数多于真实的头，说明释放机制是调谐于某一特征的过滤器，
如红色对比、细长形状，而不是对亲代的识别；
夸大该特征，
反应也随之增强。
\end{solution}

\begin{solution}{exo:b3:behavioural-ecology:3}
当斑块中的收益速率降至整个栖息地的平均速率时便离开，
此处平均速率包括路途时间。路途时间加倍会降低该平均速率，
因此觅食者在每个斑块中停留更久，
对每个斑块的消耗也更彻底，最终总体速率更低。
\end{solution}

\begin{solution}{exo:b3:behavioural-ecology:4}
进化稳定策略（ESS）是指：一旦几乎所有个体都采用它，
任何稀有的替代策略都无法获得更高收益。全鸽策略并非如此：一只单独的鹰
赢得每次争斗，对鸽子得到 $V$，而鸽子彼此之间只得到 $V/2$。在混合策略中，各策略
的收益为 $V(1 - V/C)/2$，低于全鸽时的 $V/2$，
因为有一部分争斗会造成伤害；单独的个体无法靠自身改进，
因此所有人都陷于这一损失之中。
\end{solution}

\begin{solution}{exo:b3:behavioural-ecology:5}
$V = 6$，$C = 10$：$p^{*} = 0.6$；每个个体的收益为 $6(1 - 0.6)/2 = 1.2$，
而全鸽策略的收益为 $3$。当 $V = 12 > C$ 时，纯鹰策略是进化稳定策略，其收益
为 $(12 - 10)/2 = 1$，而全鸽策略为 $6$——争斗
摧毁了六分之五的价值。
\end{solution}

\begin{solution}{exo:b3:behavioural-ecology:6}
$\tau = T_{0}$：$t^{*} = 1.15\times 2 = \qty{2.3}{min}$，比例 $1 -
\mathrm{e}^{-1.15} = 0.68$，$R = 0.68A/4.3 = 0.16A$ 每分钟。$\tau =
4T_{0}$：$t^{*} = 2.0\times 2 = \qty{4}{min}$（更精确地为 $3.9$），
比例 $0.86$，$R = 0.86A/12 = 0.072A$ 每分钟。
\end{solution}

\begin{solution}{exo:b3:behavioural-ecology:7}
同父异母（半同胞）为 $1/4$；祖父母–孙辈为 $1/4$；堂表兄弟姐妹为 $1/8$。
单双倍体工蜂与其兄弟：兄弟为单倍体，只有
其母亲的基因；工蜂的基因有一半来自母亲，
其中一半是他所得的拷贝，因此从她的角度看 $r = 1/4$，
（从他的角度为 $1/2$：倍性不同时，亲缘关系并不
对称）。蜂后与其子：他的基因全部来自她，但她的基因只有一半
在他体内：从她的一方看 $r = 1/2$，从他的一方看为 $1$。
\end{solution}

\begin{solution}{exo:b3:behavioural-ecology:8}
在 $n\,r\,b > c$ 中，设 $b = 0.01$，$c = 0.02$，则 $n\,r > 2$。全同胞，$r
= 1/2$：$n \ge 5$。侄女、外甥女，$r = 1/4$：$n \ge 9$。无亲缘关系：永不成立——
这正是扩散的雄性不鸣叫的原因。
\end{solution}

\begin{solution}{exo:b3:behavioural-ecology:9}
每 \qty{25}{cm} 交配次数加倍，存活率下降五分之一：$L = 25$：
$0.5\times 1.2 = 0.6$；$50$：$1$；$75$：$2\times 0.8 = 1.6$；$100$：
$4\times 0.6 = 2.4$；$125$：$8\times 0.4 = 3.2$；$150$：$16\times 0.2 =
3.2$。按这些条件，乘积会持续上升，远超自然的
\qty{50}{cm}，因此模型遗漏了成本：飞行中长尾的存活代价
比线性更陡，雌性的反应趋于饱和，长羽毛耗费能量，而且
安德森关于巢数加倍的测定只在一个季节内完成，
而非一生。自然的长尾
正是真实成本开始起作用之处。
\end{solution}

\begin{solution}{exo:b3:behavioural-ecology:10}
报复者：对鹰的收益为 $(V - C)/2$，对鸽子为 $V/2$，对报复者为
$V/2$。在全为报复者的种群中，稀有的鹰遇到的是会还击的报复者，
其收益为 $(V - C)/2 < V/2$：当 $V < C$ 时，它无法入侵。
稀有的鸽子收益为 $V/2$，恰与报复者相同，因此鸽子是中性的，
可以渐渐进入种群；一旦数量增多，
便会让鹰进入（鹰对鸽子的收益为 $V$），三策略博弈
没有简单的进化稳定策略——这是首个例证，说明增加一个选项
如何能够打破一个稳定的混合状态。
\end{solution}

\begin{solution}{exo:b3:behavioural-ecology:11}
若所有个体都恰好持续 $t_{0}$，则持续 $t_{0} +
\varepsilon$ 的突变体将以可忽略的额外代价赢得每次争斗，而由这种突变体组成的种群
将被持续 $t_{0} + 2\varepsilon$ 的个体击败，如此类推：
不存在稳定的固定时长，因此进化稳定策略必须是随机的。在指数分布下，
再持续 $s$ 的概率，
在已持续 $t$ 的条件下，即 $\mathrm{e}^{-s/m}$，无论 $t$ 为何值，
即无记忆性质——因此仍在展示的对手，不会透露它还会持续多久，
而任何策略都无法
利用已逝去的时间。
\end{solution}

\begin{solution}{exo:b3:behavioural-ecology:12}
工蜂对姐妹的价值估计为 $3/4$，对兄弟为 $1/4$，因此在工蜂控制下，
蜂群应在雌性上的投入为雄性的三倍，
即按投入计为 $3{:}1$；蜂后对两者亲缘关系相同，
偏好 $1{:}1$。特里弗斯和黑尔（1976）发现，
在一只单交配、仅交配一次的蜂后所产的蚂蚁中，投入比接近 $3{:}1$：工蜂获胜。
多次交配会降低工蜂与姐妹间的平均亲缘关系，
使其趋近 $1/4$ 至 $1/2$，并使工蜂的最优比例趋向 $1{:}1$；
多蜂后交配的物种的蜂群，确实投入得更为均匀，
与该论证的预测一致。
\end{solution}

\begin{solution}{pb:b3:behavioural-ecology:1}
\textbf{1.} $R(t) = A(1 - \mathrm{e}^{-t/T_{0}})/(\tau + t)$；离开的条件为
$g'(t) = R(t)$：$(A/T_{0})\mathrm{e}^{-t/T_{0}} = A(1 - \mathrm{e}^{-t/T_{0}})/
(\tau + t)$。
\textbf{2.} 两边乘以 $(\tau + t)\mathrm{e}^{t/T_{0}}/A$：$(\tau + t)/T_{0}
= \mathrm{e}^{t/T_{0}} - 1$。取 $\tau = T_{0}$，$x = t/T_{0}$，得 $2 + x =
\mathrm{e}^{x}$；$x = 1.15$ 时左边为 $3.15$，右边为 $3.16$。
\textbf{3.} $\tau = 4T_{0}$：$5 + x = \mathrm{e}^{x}$；$x = 1.94$ 时
左边为 $6.94$，右边为 $6.96$；$t^{*} \approx 1.94\,T_{0}$，取 $2.0$。
\textbf{4.} 密集斑块：$t^{*} = \qty{3.5}{min}$，$60(1 - \mathrm{e}^{-1.15}) =
\qty{41}{kJ}$，\qty{68}{\%}，$R = 41/6.5 = \qty{6.4}{kJ/min}$。稀疏斑块：
$t^{*} = \qty{5.8}{min}$，\qty{51}{kJ}，\qty{86}{\%}，$R = 51/17.8 =
\qty{2.9}{kJ/min}$。
\textbf{5.} 五分钟时：$g(5) = 60(1 - \mathrm{e}^{-5/3}) = \qty{49}{kJ}$。
密集斑块：$49/8 = \qty{6.1}{kJ/min}$（最优值的 \qty{96}{\%}）；稀疏斑块：
$49/17 = \qty{2.9}{kJ/min}$（\qty{99}{\%}）。粗略的规则会损失
百分之几的收益——对选择而言，这已经足够好，因而被采用。
\textbf{6.} 随着斑块被耗尽，两次捕获之间的间隔越来越长，
因此固定的放弃时间即对应固定的边际捕获速率——
正是定理的标准，它在每个斑块中都相同。该规则并不
可靠，因为正确的阈值取决于生境：若固定为
\qty{20}{s}，则在移动距离较长的地方动物离开得过早，而在移动距离较短的地方离开得过晚，
除非放弃时间本身能根据近期经验加以调整——
觅食者确实如此。
\textbf{7.} 鹰对鹰 $(10 - 30)/2 = -10$，鹰对鸽 $10$，鸽对
鹰 $0$，鸽对鸽 $5$；$p^{*} = V/C = 1/3$。
\textbf{8.} $W_{H} = 10 - 20p$，$W_{D} = 5(1 - p)$。$p = 0.1$：$8$ 对
$4.5$，鹰扩散。$p = 1/3$：各为 $3.33$，达到平衡。$p = 0.6$：$-2$
对 $2$，鸽扩散。
\textbf{9.} \omterm{prop:b3:behavioural-ecology:hawk-dove}{演化稳定策略}下为 $3.33$，而全部为鸽时为 $5$：每次争斗的价值中，
有三分之一因可能发生打斗而损失。
\textbf{10.} $V = 40 > C$：鹰是纯粹的\omterm{prop:b3:behavioural-ecology:hawk-dove}{演化稳定策略}；每次争斗都是一场打斗。
困难的年份应当出现更多的打斗和更多的伤害——实际情况也确实如此。
\textbf{11.} 鹰遇到资产者时，有一半的时间它是占有者
（打斗，收益 $(V - C)/2$），另一半时间它是入侵者（不战而胜，收益 $V$）：
平均收益为 $(3V - C)/4$。资产者在与自身同类相遇时收益为 $V/2$。仅当
$(3V - C)/4 > V/2$，即 $V > C$ 时，鹰才能入侵。蝴蝶的“占有者
获胜”即为资产者策略：这是一种代价低廉的约定，之所以稳定，是因为打斗的代价
超过了一块太阳黑子的价值——而当双方都自认为是占有者时，
双方都会采取鹰的策略。
\textbf{12.} 最佳策略取决于其他所有个体的行为，
因此种群会稳定在各策略收益相等
的状态，而不是任何个体的最佳状态。
\textbf{13.} $n\bar r\,b > c$：$n\bar r > c/b = 3$。
\textbf{14.} 八只全兄妹：$n\bar r = 4 > 3$，值班有利可图。八只
陌生个体：$0$，从不值班。八只同父异母或同母异父的兄弟姐妹：$2 < 3$，不值班。
\textbf{15.} $c = 0.01$：阈值为 $n\bar r > 1$——此时同父异母或同母异父的兄弟姐妹，
甚至少数全兄妹，都已足够。成本最小的个体
恰好满足该规则，因此吃饱的个体站岗放哨，
而饥饿的个体去挖洞。
\textbf{16.} $k$ 个姐妹的价值为 $\tfrac{3}{4}k$，$k$ 个女儿
的价值为 $\tfrac{1}{2}k$：偏好姐妹，比例为 $3{:}2$。
\textbf{17.} 随机一个姐妹是全姐妹（$3/4$）的概率为 $1/3$，是同父异母或同母异父的
姐妹（$1/4$）的概率为 $2/3$：$\bar r = 0.25 +
0.17 = 0.42 < 0.5$。偏好发生逆转：此时女儿胜过姐妹，
因此单靠亲缘关系无法解释一妻多夫物种中的不育工蜂——
还须加入生态因素与蜂群层面的监管。
\textbf{18.} 该规则只要求被帮助者的 $r$ 平均高于种群平均水平；
任何与亲缘关系相关的线索都可以起作用——
留在出生洞穴附近，帮助与自己一同长大的个体。
地松鼠的规则可能是“在离家近时鸣叫”，它借此筛选出亲属，
却无需识别任何个体。
\textbf{19.} $W = (L/20)(1 - L^{2}/10^{4})$：$L = 20$：$0.96$；$40$：
$1.68$；$60$：$1.92$；$80$：$1.44$。
\textbf{20.} $W' = \tfrac{1}{20}(1 - 3L^{2}/10^{4}) = 0$：$L^{*} = 100/\sqrt{3}
= \qty{58}{cm}$，$W = 2.89\times\tfrac{2}{3} = 1.92$。
\textbf{21.} $L^{*} = 70/\sqrt{3} = \qty{40}{cm}$：更短。最优点位于交配边际收益
等于生存边际损失之处，因此捕食成本越高，尾羽就越短；
处于不同捕食压力下的种群，
其装饰性状应当有所不同——孔雀鱼与
剑尾鱼确实如此。
\textbf{22.} 两者皆是。\omterm{def:b3:behavioural-ecology:sexual-selection}{失控选择}预测偏好会超过性状本身，
而性状止步于生存最优值；累赘假说则预测雌性偏好更长的尾羽，
因为只有强健的雄性才能负担得起，这又超出了大多数个体所能生长的长度。
该实验显示了超过最优值的偏好，
但并不能区分两种解释。
\textbf{23.} 若尾羽无成本，每只雄性都能长出最长的尾羽，
它便无法反映其质量，雌性对其的偏好也会被利用而逐渐消退。
若成本对状况较差的雄性负担更重，
则只有强健的个体才承担得起长尾，
此时偏好便会选择适合度：成本正是使
信号可信的原因。
\textbf{24.} 每次交配的回报为八倍，因此雄性的适合度随交配次数急剧上升，
而雌性几乎不随之增加；雄性受到选择，
需争夺交配机会——即鹰–鸽争斗及其武器——
而雌性的后代数量虽少，但每一个都很重要，因而受到选择去挑选配偶，
这使得雄性的装饰性状有利可图。
每次交配收益上的一种不对称，同时产生了争斗与华丽的装饰。
\textbf{25.} 密集斑块中 $t^{*} \approx \qty{3.5}{min}$；$p^{*}
= 1/3$；当 $n\bar r > 3$ 时，站岗放哨有利可图。
\end{solution}
